\section{Preliminaries}\label{prelim}
We start with some useful properties of the quasi-invariants.
\begin{Proposition}[\cite{etingof2002lectures}]\label{firstProp}
 Let $\k$ be a field.
\begin{enumerate}\itemsep=0pt
 \item[$1$.] $\k[x_1,x_2,x_3]^{S_3}\subset Q_m(3,\k)$, $Q_0(3,\k) = \k[x_1,x_2,x_3]$, and $Q_m(3,\k)\supset Q_{m'}(3,\k)$, where ${m'>m}$.
 \item[$2$.] $Q_m(3,\k)$ is a ring.
 \item[$3$.] $Q_m(3,\k)$ is a finitely generated $\k[x_1,x_2,x_3]^{S_3}$-module.
\end{enumerate}
\end{Proposition}

Note that \cite{etingof2002lectures} proves Proposition~\ref{firstProp} in the case, where $\k=\C$. However, the proofs for the first two assertions work over any field, and the last assertion follows from the Hilbert basis theorem. In view of the structure of $Q_m(3,\k)$ as a module over the symmetric polynomials, given some $K\in Q_m(3,\k)$, we will frequently refer to quasi-invariant polynomials that can be obtained via scalar multiplication of $Q$ by a symmetric polynomial. To distinguish these polynomials from the ordinary $\k$-multiples of $Q$, we will refer to them as \emph{symmetric polynomial multiples} of $Q$.

We consider $Q_m(3, \mathbf{F}_{2})$ and $Q_m(3, \mathbf{F}_{3})$ as representations of $S_3$, where $S_3$ permutes the variables $x_1$, $x_2$, $x_3$. Since $Q_m(3, \mathbf{F}_{2})$ and $Q_m(3, \mathbf{F}_{3})$ are vector spaces over $\mathbf{F}_2$ and $\mathbf{F}_3$ respectively and the characteristics $2$ and $3$ divide $|S_3|$, $Q_m(3, \mathbf{F}_{2})$ and $Q_m(3, \mathbf{F}_{3})$ are modular representations~of~$S_3$.

\begin{Proposition}
 $Q_m(3, \mathbf{F}_{2})$ and $Q_m(3, \mathbf{F}_{3})$ are modular representations of $S_3$.
\end{Proposition}

First, we consider characteristic $2$.

